% titel: Blatt 9, A1: Vollständige Induktion
% Teil II: der durchgerechnete Weg.

\subsection*{Gegeben und gesucht}
Zwei Aussagen für alle $n\in\mathbb N$ per vollständiger Induktion zeigen:
\begin{align}
  \text{a)}\quad &\sum_{k=1}^{n}\frac{k}{2^k}=2-\frac{n+2}{2^n}, \\
  \text{b)}\quad &\prod_{k=1}^{n}\Bigl(1+\frac2k\Bigr)=\frac{(n+1)(n+2)}2.
\end{align}

\subsection*{a) Summe $\sum_{k=1}^{n}\frac{k}{2^k}$}
Induktionsanfang ($n=1$): $\sum_{k=1}^{1}\frac{k}{2^k}=\frac12$ und
$2-\frac{3}{2}=\frac12$ — stimmt.

Induktionsschritt ($n\to n+1$): mit der Induktionsannahme
$\sum_{k=1}^{n}\frac{k}{2^k}=2-\frac{n+2}{2^n}$ gilt
\begin{align}
  \sum_{k=1}^{n+1}\frac{k}{2^k}
  &=\sum_{k=1}^{n}\frac{k}{2^k}+\frac{n+1}{2^{n+1}} \\
  &=2-\frac{n+2}{2^n}+\frac{n+1}{2^{n+1}}
  =2-\frac{2(n+2)-(n+1)}{2^{n+1}} \\
  &=2-\frac{n+3}{2^{n+1}}
  =2-\frac{(n+1)+2}{2^{n+1}}.
\end{align}
Das ist die Behauptung für $n+1$. Aus IA und IS folgt die Formel für alle
$n\in\mathbb N$.

\begin{ergebnis}
$\displaystyle\sum_{k=1}^{n}\frac{k}{2^k}=2-\frac{n+2}{2^n}\quad\forall n\in\mathbb N$.
\end{ergebnis}

\subsection*{b) Produkt $\prod_{k=1}^{n}\bigl(1+\tfrac2k\bigr)$}
Schreibe $1+\frac2k=\frac{k+2}{k}$. Der Bruch kürzt jedes Mal die vorige
Zahl:
\begin{align}
  \prod_{k=1}^{n}\frac{k+2}{k}
  &=\frac{3}{1}\cdot\frac{4}{2}\cdot\frac{5}{3}\cdots\frac{n+1}{n-1}\cdot\frac{n+2}{n} \\
  &=\frac{(n+1)\,(n+2)}{1\cdot 2}
  =\frac{(n+1)(n+2)}2.
\end{align}
Alternativ sauber induktiv: IA ($n=1$) gibt $\frac{3}{1}\cdot\frac{4}{2}\cdot\frac{5}{3}=\frac{3}{1}=3$ links und $\frac{2\cdot3}{2}=3$ rechts. Im IS multipliziert man mit
dem Faktor $1+\frac{2}{n+1}=\frac{n+3}{n+1}$ und erhält
$\frac{(n+1)(n+2)}{2}\cdot\frac{n+3}{n+1}=\frac{(n+2)(n+3)}2$, die
Behauptung für $n+1$.

\begin{ergebnis}
$\displaystyle\prod_{k=1}^{n}\Bigl(1+\frac2k\Bigr)=\frac{(n+1)(n+2)}2\quad\forall n\in\mathbb N$.
\end{ergebnis}

\begin{kontrolle}
\textbf{Kleiner Fall:} für $n=1$ ist im Produkt nur der Faktor $1+\frac21=3$
links, rechts $\frac{2\cdot3}{2}=3$ — passt. \textbf{Grenzfall:} für große
$n$ wächst das Produkt wie $\frac{n^2}{2}$, der Ausdruck rechts ebenso —
plausibel.
\end{kontrolle}

\begin{merksatz}
Induktion heißt: IA prüfen, dann den neuen Summanden/Faktor abspalten und die
Voraussetzung einsetzen; Ausdrücke wie $\frac{k+2}{k}$ als Teleskopprodukt zu
schreiben vereinfacht den Rückweg auf die Annahme.
\end{merksatz}
